\section{账本事件的社会制度深描：wanli}

\label{sec:ledger-depth-wanli}

以下各节依本卷账本的事件、来源和置信提示组织，不以编年节点替代地方社会。朝廷命令、法律条文、城市记录、家庭文书与物质遗存的证明范围不同，必须让其差异保留在叙事中。

\subsection{政治制度与地方执行}

在账本条目\texttt{undefined}中，1572（隆庆六年）六月的“高拱被罢，张居正成为万历初年内阁的核心人物”发生于明。本条把背景说明为幼主继位后，内阁首辅、太后与司礼监围绕诏令传达和用人权发生冲突。，并以张居正获得推行考成、清丈和赋役改革的政治空间；内阁与司礼监合作也成为其权力基础和争议来源。概括可见影响。要理解它的社会后果，不能止于宫廷或战场：土地、赋役、司法、道路、性别化家庭劳动、城市市场、书院寺观和边海网络都会影响地方如何接收或改写制度。

本条依据《穆宗实录》隆庆六年六月条；《明史·高拱传》；《明史·张居正传》。这些来源能够较稳妥地支持所记年序、诏令、任命或特定地点的行动，但无法自动证明全国的财产、人口、识字、信仰或动机。账本保留的证据说明是“罢黜和张居正地位上升可靠；关于太后、冯保和张居正的密谋细节多来自后出叙事，难以独立验证。”。所以，官府标签、奏疏数字与后来的道德评语必须放回其文书目的，不能被写成不经分区核验的社会全貌。

这条事件更适合作为一组关系变化的锚点：中央方案需由地方官、书吏、士绅、宗族、商人与劳动者共同转译；不同家庭面对的税役、婚姻、迁徙和安全风险并不相等。材料只让我们看见其中有限的记录者。承认可见与不可见的界线，才能使制度史、文化史和区域史互相校正。

在账本条目\texttt{undefined}中，1572（隆庆六年）五月的“隆庆帝去世，十岁的朱翊钧继位，是为万历帝”发生于明。本条把背景说明为朱翊钧已为皇太子，但幼主即位使太后、司礼监和内阁在政务传达中具有更大作用。，并以高拱与张居正围绕辅政位置的竞争迅速公开化，万历初年改革的政治条件由此形成。概括可见影响。要理解它的社会后果，不能止于宫廷或战场：土地、赋役、司法、道路、性别化家庭劳动、城市市场、书院寺观和边海网络都会影响地方如何接收或改写制度。

本条依据《穆宗实录》隆庆六年五月条；《神宗实录》万历元年即位条；《明史·神宗本纪》。这些来源能够较稳妥地支持所记年序、诏令、任命或特定地点的行动，但无法自动证明全国的财产、人口、识字、信仰或动机。账本保留的证据说明是“死亡、继位与嗣位年龄可靠；太后、司礼监和内阁的实际分工须以具体诏令、批答和人事记录判断。”。所以，官府标签、奏疏数字与后来的道德评语必须放回其文书目的，不能被写成不经分区核验的社会全貌。

这条事件更适合作为一组关系变化的锚点：中央方案需由地方官、书吏、士绅、宗族、商人与劳动者共同转译；不同家庭面对的税役、婚姻、迁徙和安全风险并不相等。材料只让我们看见其中有限的记录者。承认可见与不可见的界线，才能使制度史、文化史和区域史互相校正。

在账本条目\texttt{undefined}中，1573（万历元年）的“西班牙与中国的贸易从月港开始”发生于明与海上、朝贡相关政权。本条把背景说明为朝贡名义、边境安全与港口贸易的利益使交涉持续发生。，并以它为后续册封、互市或海防安排留下文书与跨政权比较线索。概括可见影响。要理解它的社会后果，不能止于宫廷或战场：土地、赋役、司法、道路、性别化家庭劳动、城市市场、书院寺观和边海网络都会影响地方如何接收或改写制度。

本条依据《神宗实录》万历元年条；《明史·外国传》；《朝鲜王朝实录》、琉球《历代宝案》或海外档案。这些来源能够较稳妥地支持所记年序、诏令、任命或特定地点的行动，但无法自动证明全国的财产、人口、识字、信仰或动机。账本保留的证据说明是“译写候选以编年材料定位；实录记载反映朝廷选择，崇祯期须另以奏疏、地方及敌对政权材料交叉。”。所以，官府标签、奏疏数字与后来的道德评语必须放回其文书目的，不能被写成不经分区核验的社会全貌。

这条事件更适合作为一组关系变化的锚点：中央方案需由地方官、书吏、士绅、宗族、商人与劳动者共同转译；不同家庭面对的税役、婚姻、迁徙和安全风险并不相等。材料只让我们看见其中有限的记录者。承认可见与不可见的界线，才能使制度史、文化史和区域史互相校正。

在账本条目\texttt{undefined}中，1573（万历元年）的“俺答边患、东南海防或沿海武装的本年军令与战报构成危机处理线索，战果和参与者须交叉核对。”发生于明。本条把背景说明为前期边防、地方武装或跨区域资源调动的压力推动了该行动。，并以它改变了后续调兵、补给或地方秩序的条件；战果和伤亡须分开核对。概括可见影响。要理解它的社会后果，不能止于宫廷或战场：土地、赋役、司法、道路、性别化家庭劳动、城市市场、书院寺观和边海网络都会影响地方如何接收或改写制度。

本条依据《神宗实录》万历元年条；《明史·兵志》及相关列传；边镇、地方志或对方材料。这些来源能够较稳妥地支持所记年序、诏令、任命或特定地点的行动，但无法自动证明全国的财产、人口、识字、信仰或动机。账本保留的证据说明是“桥接条目只标示可回查的年度问题与材料链；实录有中央选择性，崇祯期尤其需要奏疏、地方、朝鲜及满文材料交叉。”。所以，官府标签、奏疏数字与后来的道德评语必须放回其文书目的，不能被写成不经分区核验的社会全貌。

这条事件更适合作为一组关系变化的锚点：中央方案需由地方官、书吏、士绅、宗族、商人与劳动者共同转译；不同家庭面对的税役、婚姻、迁徙和安全风险并不相等。材料只让我们看见其中有限的记录者。承认可见与不可见的界线，才能使制度史、文化史和区域史互相校正。

在账本条目\texttt{undefined}中，1573（万历元年）的“严嵩时期至隆庆初的人事和奏议材料标记中枢权力变化，案狱与党争标签需保留来源立场。”发生于明。本条把背景说明为继承、官僚协调或皇权运作的压力使问题进入朝廷程序。，并以它影响后续人事与政策议程；动机和派系标签不能由单一叙事裁定。概括可见影响。要理解它的社会后果，不能止于宫廷或战场：土地、赋役、司法、道路、性别化家庭劳动、城市市场、书院寺观和边海网络都会影响地方如何接收或改写制度。

本条依据《神宗实录》万历元年条；《明史·本纪》及相关列传；诏令、奏疏或会典版本。这些来源能够较稳妥地支持所记年序、诏令、任命或特定地点的行动，但无法自动证明全国的财产、人口、识字、信仰或动机。账本保留的证据说明是“桥接条目只标示可回查的年度问题与材料链；实录有中央选择性，崇祯期尤其需要奏疏、地方、朝鲜及满文材料交叉。”。所以，官府标签、奏疏数字与后来的道德评语必须放回其文书目的，不能被写成不经分区核验的社会全貌。

这条事件更适合作为一组关系变化的锚点：中央方案需由地方官、书吏、士绅、宗族、商人与劳动者共同转译；不同家庭面对的税役、婚姻、迁徙和安全风险并不相等。材料只让我们看见其中有限的记录者。承认可见与不可见的界线，才能使制度史、文化史和区域史互相校正。

在账本条目\texttt{undefined}中，1573（万历元年）的“张居正推行考成法，要求部院与地方官限期申报、完成政务”发生于明。本条把背景说明为万历初年中央希望缩短命令滞留并掌握各部、各省的执行进度，张居正获得内阁与司礼监支持。，并以考成加强了限期考核和责任追究，也促使官员以文书合规回应压力；其效果须与地方事务逐项区分。概括可见影响。要理解它的社会后果，不能止于宫廷或战场：土地、赋役、司法、道路、性别化家庭劳动、城市市场、书院寺观和边海网络都会影响地方如何接收或改写制度。

本条依据《神宗实录》万历元年吏部条；《明史·张居正传》；张居正奏疏与地方档案。这些来源能够较稳妥地支持所记年序、诏令、任命或特定地点的行动，但无法自动证明全国的财产、人口、识字、信仰或动机。账本保留的证据说明是“考成规定和张居正角色可靠；报送完成率反映行政压力与文书策略，不等同政策在各地真实落实。”。所以，官府标签、奏疏数字与后来的道德评语必须放回其文书目的，不能被写成不经分区核验的社会全貌。

这条事件更适合作为一组关系变化的锚点：中央方案需由地方官、书吏、士绅、宗族、商人与劳动者共同转译；不同家庭面对的税役、婚姻、迁徙和安全风险并不相等。材料只让我们看见其中有限的记录者。承认可见与不可见的界线，才能使制度史、文化史和区域史互相校正。

在账本条目\texttt{undefined}中，1573（万历元年）的“李成梁进攻王杲集团，王杲被擒后由明廷处置”发生于明与建州女真、辽东。本条把背景说明为王杲与辽东边市、朝贡及部落竞争相连，李成梁试图以军事和招抚重整边境关系。，并以王杲集团受挫改变了建州内部力量平衡；努尔哈赤早年的政治处境须避免从后来的后金兴起倒推。概括可见影响。要理解它的社会后果，不能止于宫廷或战场：土地、赋役、司法、道路、性别化家庭劳动、城市市场、书院寺观和边海网络都会影响地方如何接收或改写制度。

本条依据《神宗实录》万历元年辽东条；《明史·李成梁传》；《清实录》前代建州相关追述。这些来源能够较稳妥地支持所记年序、诏令、任命或特定地点的行动，但无法自动证明全国的财产、人口、识字、信仰或动机。账本保留的证据说明是“进攻、被擒和朝廷处置可靠；女真各部关系与李成梁战功数字不能仅按辽东奏报或后世“养成建州”叙事简化。”。所以，官府标签、奏疏数字与后来的道德评语必须放回其文书目的，不能被写成不经分区核验的社会全貌。

这条事件更适合作为一组关系变化的锚点：中央方案需由地方官、书吏、士绅、宗族、商人与劳动者共同转译；不同家庭面对的税役、婚姻、迁徙和安全风险并不相等。材料只让我们看见其中有限的记录者。承认可见与不可见的界线，才能使制度史、文化史和区域史互相校正。

在账本条目\texttt{undefined}中，1574（万历二年）的“李成梁在乔仓嘎和塔克西的帮助下杀死了王高”发生于明。本条把背景说明为继承、官僚协调或皇权运作的压力使问题进入朝廷程序。，并以它影响后续人事与政策议程；动机和派系标签不能由单一叙事裁定。概括可见影响。要理解它的社会后果，不能止于宫廷或战场：土地、赋役、司法、道路、性别化家庭劳动、城市市场、书院寺观和边海网络都会影响地方如何接收或改写制度。

本条依据《神宗实录》万历二年条；《明史·本纪》及相关列传；诏令、奏疏或会典版本。这些来源能够较稳妥地支持所记年序、诏令、任命或特定地点的行动，但无法自动证明全国的财产、人口、识字、信仰或动机。账本保留的证据说明是“译写候选以编年材料定位；实录记载反映朝廷选择，崇祯期须另以奏疏、地方及敌对政权材料交叉。”。所以，官府标签、奏疏数字与后来的道德评语必须放回其文书目的，不能被写成不经分区核验的社会全貌。

这条事件更适合作为一组关系变化的锚点：中央方案需由地方官、书吏、士绅、宗族、商人与劳动者共同转译；不同家庭面对的税役、婚姻、迁徙和安全风险并不相等。材料只让我们看见其中有限的记录者。承认可见与不可见的界线，才能使制度史、文化史和区域史互相校正。

在账本条目\texttt{undefined}中，1574（万历二年）的“澳门周围竖起了围墙”发生于明。本条把背景说明为继承、官僚协调或皇权运作的压力使问题进入朝廷程序。，并以它影响后续人事与政策议程；动机和派系标签不能由单一叙事裁定。概括可见影响。要理解它的社会后果，不能止于宫廷或战场：土地、赋役、司法、道路、性别化家庭劳动、城市市场、书院寺观和边海网络都会影响地方如何接收或改写制度。

本条依据《神宗实录》万历二年条；《明史·本纪》及相关列传；诏令、奏疏或会典版本。这些来源能够较稳妥地支持所记年序、诏令、任命或特定地点的行动，但无法自动证明全国的财产、人口、识字、信仰或动机。账本保留的证据说明是“译写候选以编年材料定位；实录记载反映朝廷选择，崇祯期须另以奏疏、地方及敌对政权材料交叉。”。所以，官府标签、奏疏数字与后来的道德评语必须放回其文书目的，不能被写成不经分区核验的社会全貌。

这条事件更适合作为一组关系变化的锚点：中央方案需由地方官、书吏、士绅、宗族、商人与劳动者共同转译；不同家庭面对的税役、婚姻、迁徙和安全风险并不相等。材料只让我们看见其中有限的记录者。承认可见与不可见的界线，才能使制度史、文化史和区域史互相校正。

在账本条目\texttt{undefined}中，1574（万历二年）的“本年灾情、赈恤或河工记录连接中央与地方社会，区域差异需以府县材料追踪。”发生于明。本条把背景说明为自然风险与地方报告促成朝廷或地方的应对记录。，并以赈济、迁徙和损失的实际范围仍须以受灾地区材料分别核验。概括可见影响。要理解它的社会后果，不能止于宫廷或战场：土地、赋役、司法、道路、性别化家庭劳动、城市市场、书院寺观和边海网络都会影响地方如何接收或改写制度。

本条依据《神宗实录》万历二年条；《明史·五行志》；受灾府县地方志与奏疏。这些来源能够较稳妥地支持所记年序、诏令、任命或特定地点的行动，但无法自动证明全国的财产、人口、识字、信仰或动机。账本保留的证据说明是“桥接条目只标示可回查的年度问题与材料链；实录有中央选择性，崇祯期尤其需要奏疏、地方、朝鲜及满文材料交叉。”。所以，官府标签、奏疏数字与后来的道德评语必须放回其文书目的，不能被写成不经分区核验的社会全貌。

这条事件更适合作为一组关系变化的锚点：中央方案需由地方官、书吏、士绅、宗族、商人与劳动者共同转译；不同家庭面对的税役、婚姻、迁徙和安全风险并不相等。材料只让我们看见其中有限的记录者。承认可见与不可见的界线，才能使制度史、文化史和区域史互相校正。

在账本条目\texttt{undefined}中，1574（万历二年）的“军饷、盐课、漕运和赋役的本年行政材料显示中晚明财政压力，法定额与实收不可互代。”发生于明。本条把背景说明为财政征解、交通工程或货币支付的压力使相关措施进入政务。，并以其法定安排不等于地方执行，后续须以地域材料追踪负担与成效。概括可见影响。要理解它的社会后果，不能止于宫廷或战场：土地、赋役、司法、道路、性别化家庭劳动、城市市场、书院寺观和边海网络都会影响地方如何接收或改写制度。

本条依据《神宗实录》万历二年条；《明会典》相关条目；地方志、黄册/契约/河工材料。这些来源能够较稳妥地支持所记年序、诏令、任命或特定地点的行动，但无法自动证明全国的财产、人口、识字、信仰或动机。账本保留的证据说明是“桥接条目只标示可回查的年度问题与材料链；实录有中央选择性，崇祯期尤其需要奏疏、地方、朝鲜及满文材料交叉。”。所以，官府标签、奏疏数字与后来的道德评语必须放回其文书目的，不能被写成不经分区核验的社会全貌。

这条事件更适合作为一组关系变化的锚点：中央方案需由地方官、书吏、士绅、宗族、商人与劳动者共同转译；不同家庭面对的税役、婚姻、迁徙和安全风险并不相等。材料只让我们看见其中有限的记录者。承认可见与不可见的界线，才能使制度史、文化史和区域史互相校正。

\subsection{法律、家庭与社会关系}

在账本条目\texttt{undefined}中，1574（万历二年）的“海防知识、学校、科举与礼制的本年文本提供制度和知识史线索，方案不等于实际配置。”发生于明。本条把背景说明为制度、教育或知识传播的需要推动了相关行动或文本形成。，并以它提供制度和传播史的锚点，不能据此推断社会整体接受。概括可见影响。要理解它的社会后果，不能止于宫廷或战场：土地、赋役、司法、道路、性别化家庭劳动、城市市场、书院寺观和边海网络都会影响地方如何接收或改写制度。

本条依据《神宗实录》万历二年条；《明史·选举志》或《艺文志》；文集、碑刻或书目材料。这些来源能够较稳妥地支持所记年序、诏令、任命或特定地点的行动，但无法自动证明全国的财产、人口、识字、信仰或动机。账本保留的证据说明是“桥接条目只标示可回查的年度问题与材料链；实录有中央选择性，崇祯期尤其需要奏疏、地方、朝鲜及满文材料交叉。”。所以，官府标签、奏疏数字与后来的道德评语必须放回其文书目的，不能被写成不经分区核验的社会全貌。

这条事件更适合作为一组关系变化的锚点：中央方案需由地方官、书吏、士绅、宗族、商人与劳动者共同转译；不同家庭面对的税役、婚姻、迁徙和安全风险并不相等。材料只让我们看见其中有限的记录者。承认可见与不可见的界线，才能使制度史、文化史和区域史互相校正。

在账本条目\texttt{undefined}中，1575（万历三年）的“明朝海军军官王旺高抵达吕宋岛，并随马丁·德·拉达率领的西班牙大使馆返回；由于西班牙无法抓获中国海盗林峰，使馆失败”发生于明与海上、朝贡相关政权。本条把背景说明为朝贡名义、边境安全与港口贸易的利益使交涉持续发生。，并以它为后续册封、互市或海防安排留下文书与跨政权比较线索。概括可见影响。要理解它的社会后果，不能止于宫廷或战场：土地、赋役、司法、道路、性别化家庭劳动、城市市场、书院寺观和边海网络都会影响地方如何接收或改写制度。

本条依据《神宗实录》万历三年条；《明史·外国传》；《朝鲜王朝实录》、琉球《历代宝案》或海外档案。这些来源能够较稳妥地支持所记年序、诏令、任命或特定地点的行动，但无法自动证明全国的财产、人口、识字、信仰或动机。账本保留的证据说明是“译写候选以编年材料定位；实录记载反映朝廷选择，崇祯期须另以奏疏、地方及敌对政权材料交叉。”。所以，官府标签、奏疏数字与后来的道德评语必须放回其文书目的，不能被写成不经分区核验的社会全貌。

这条事件更适合作为一组关系变化的锚点：中央方案需由地方官、书吏、士绅、宗族、商人与劳动者共同转译；不同家庭面对的税役、婚姻、迁徙和安全风险并不相等。材料只让我们看见其中有限的记录者。承认可见与不可见的界线，才能使制度史、文化史和区域史互相校正。

在账本条目\texttt{undefined}中，1575（万历三年）的“考成、人事、立储和留中等本年材料标记皇帝、内阁与言官的协调关系，不可只用党争标签解释。（材料核查重点：诏令或奏议的形成与发受链）”发生于明。本条把背景说明为继承、官僚协调或皇权运作的压力使问题进入朝廷程序。，并以它影响后续人事与政策议程；动机和派系标签不能由单一叙事裁定。概括可见影响。要理解它的社会后果，不能止于宫廷或战场：土地、赋役、司法、道路、性别化家庭劳动、城市市场、书院寺观和边海网络都会影响地方如何接收或改写制度。

本条依据《神宗实录》万历三年条；《明史·本纪》及相关列传；诏令、奏疏或会典版本。这些来源能够较稳妥地支持所记年序、诏令、任命或特定地点的行动，但无法自动证明全国的财产、人口、识字、信仰或动机。账本保留的证据说明是“桥接条目只标示可回查的年度问题与材料链；实录有中央选择性，崇祯期尤其需要奏疏、地方、朝鲜及满文材料交叉。；本条特别核对诏令或奏议的形成与发受链。”。所以，官府标签、奏疏数字与后来的道德评语必须放回其文书目的，不能被写成不经分区核验的社会全貌。

这条事件更适合作为一组关系变化的锚点：中央方案需由地方官、书吏、士绅、宗族、商人与劳动者共同转译；不同家庭面对的税役、婚姻、迁徙和安全风险并不相等。材料只让我们看见其中有限的记录者。承认可见与不可见的界线，才能使制度史、文化史和区域史互相校正。

在账本条目\texttt{undefined}中，1575（万历三年）的“清丈、一条鞭、漕运或军饷的本年文书构成万历财政的可核查节点，定额、报数和实收必须分列。”发生于明。本条把背景说明为财政征解、交通工程或货币支付的压力使相关措施进入政务。，并以其法定安排不等于地方执行，后续须以地域材料追踪负担与成效。概括可见影响。要理解它的社会后果，不能止于宫廷或战场：土地、赋役、司法、道路、性别化家庭劳动、城市市场、书院寺观和边海网络都会影响地方如何接收或改写制度。

本条依据《神宗实录》万历三年条；《明会典》相关条目；地方志、黄册/契约/河工材料。这些来源能够较稳妥地支持所记年序、诏令、任命或特定地点的行动，但无法自动证明全国的财产、人口、识字、信仰或动机。账本保留的证据说明是“桥接条目只标示可回查的年度问题与材料链；实录有中央选择性，崇祯期尤其需要奏疏、地方、朝鲜及满文材料交叉。”。所以，官府标签、奏疏数字与后来的道德评语必须放回其文书目的，不能被写成不经分区核验的社会全貌。

这条事件更适合作为一组关系变化的锚点：中央方案需由地方官、书吏、士绅、宗族、商人与劳动者共同转译；不同家庭面对的税役、婚姻、迁徙和安全风险并不相等。材料只让我们看见其中有限的记录者。承认可见与不可见的界线，才能使制度史、文化史和区域史互相校正。

在账本条目\texttt{undefined}中，1575（万历三年）的“考成、人事、立储和留中等本年材料标记皇帝、内阁与言官的协调关系，不可只用党争标签解释。（材料核查重点：报称信息的来源、抄传与行政分类）”发生于明。本条把背景说明为继承、官僚协调或皇权运作的压力使问题进入朝廷程序。，并以它影响后续人事与政策议程；动机和派系标签不能由单一叙事裁定。概括可见影响。要理解它的社会后果，不能止于宫廷或战场：土地、赋役、司法、道路、性别化家庭劳动、城市市场、书院寺观和边海网络都会影响地方如何接收或改写制度。

本条依据《神宗实录》万历三年条；《明史·本纪》及相关列传；诏令、奏疏或会典版本。这些来源能够较稳妥地支持所记年序、诏令、任命或特定地点的行动，但无法自动证明全国的财产、人口、识字、信仰或动机。账本保留的证据说明是“桥接条目只标示可回查的年度问题与材料链；实录有中央选择性，崇祯期尤其需要奏疏、地方、朝鲜及满文材料交叉。；本条特别核对报称信息的来源、抄传与行政分类。”。所以，官府标签、奏疏数字与后来的道德评语必须放回其文书目的，不能被写成不经分区核验的社会全貌。

这条事件更适合作为一组关系变化的锚点：中央方案需由地方官、书吏、士绅、宗族、商人与劳动者共同转译；不同家庭面对的税役、婚姻、迁徙和安全风险并不相等。材料只让我们看见其中有限的记录者。承认可见与不可见的界线，才能使制度史、文化史和区域史互相校正。

在账本条目\texttt{undefined}中，1575（万历三年）的“书院、科举、出版与宗教活动的本年线索可连接知识网络，存世文本有阶层与地域偏差。”发生于明。本条把背景说明为制度、教育或知识传播的需要推动了相关行动或文本形成。，并以它提供制度和传播史的锚点，不能据此推断社会整体接受。概括可见影响。要理解它的社会后果，不能止于宫廷或战场：土地、赋役、司法、道路、性别化家庭劳动、城市市场、书院寺观和边海网络都会影响地方如何接收或改写制度。

本条依据《神宗实录》万历三年条；《明史·选举志》或《艺文志》；文集、碑刻或书目材料。这些来源能够较稳妥地支持所记年序、诏令、任命或特定地点的行动，但无法自动证明全国的财产、人口、识字、信仰或动机。账本保留的证据说明是“桥接条目只标示可回查的年度问题与材料链；实录有中央选择性，崇祯期尤其需要奏疏、地方、朝鲜及满文材料交叉。”。所以，官府标签、奏疏数字与后来的道德评语必须放回其文书目的，不能被写成不经分区核验的社会全貌。

这条事件更适合作为一组关系变化的锚点：中央方案需由地方官、书吏、士绅、宗族、商人与劳动者共同转译；不同家庭面对的税役、婚姻、迁徙和安全风险并不相等。材料只让我们看见其中有限的记录者。承认可见与不可见的界线，才能使制度史、文化史和区域史互相校正。

在账本条目\texttt{undefined}中，1575（万历三年）的“辽东、西北、朝鲜或西南军务的本年调兵与战报记录跨区域动员，补给与兵额需要独立核验。”发生于明。本条把背景说明为前期边防、地方武装或跨区域资源调动的压力推动了该行动。，并以它改变了后续调兵、补给或地方秩序的条件；战果和伤亡须分开核对。概括可见影响。要理解它的社会后果，不能止于宫廷或战场：土地、赋役、司法、道路、性别化家庭劳动、城市市场、书院寺观和边海网络都会影响地方如何接收或改写制度。

本条依据《神宗实录》万历三年条；《明史·兵志》及相关列传；边镇、地方志或对方材料。这些来源能够较稳妥地支持所记年序、诏令、任命或特定地点的行动，但无法自动证明全国的财产、人口、识字、信仰或动机。账本保留的证据说明是“桥接条目只标示可回查的年度问题与材料链；实录有中央选择性，崇祯期尤其需要奏疏、地方、朝鲜及满文材料交叉。”。所以，官府标签、奏疏数字与后来的道德评语必须放回其文书目的，不能被写成不经分区核验的社会全貌。

这条事件更适合作为一组关系变化的锚点：中央方案需由地方官、书吏、士绅、宗族、商人与劳动者共同转译；不同家庭面对的税役、婚姻、迁徙和安全风险并不相等。材料只让我们看见其中有限的记录者。承认可见与不可见的界线，才能使制度史、文化史和区域史互相校正。

在账本条目\texttt{undefined}中，1576（万历四年）的“中美贸易建立”发生于明。本条把背景说明为财政征解、交通工程或货币支付的压力使相关措施进入政务。，并以其法定安排不等于地方执行，后续须以地域材料追踪负担与成效。概括可见影响。要理解它的社会后果，不能止于宫廷或战场：土地、赋役、司法、道路、性别化家庭劳动、城市市场、书院寺观和边海网络都会影响地方如何接收或改写制度。

本条依据《神宗实录》万历四年条；《明会典》相关条目；地方志、黄册/契约/河工材料。这些来源能够较稳妥地支持所记年序、诏令、任命或特定地点的行动，但无法自动证明全国的财产、人口、识字、信仰或动机。账本保留的证据说明是“译写候选以编年材料定位；实录记载反映朝廷选择，崇祯期须另以奏疏、地方及敌对政权材料交叉。”。所以，官府标签、奏疏数字与后来的道德评语必须放回其文书目的，不能被写成不经分区核验的社会全貌。

这条事件更适合作为一组关系变化的锚点：中央方案需由地方官、书吏、士绅、宗族、商人与劳动者共同转译；不同家庭面对的税役、婚姻、迁徙和安全风险并不相等。材料只让我们看见其中有限的记录者。承认可见与不可见的界线，才能使制度史、文化史和区域史互相校正。

在账本条目\texttt{undefined}中，1576（万历四年）的“辽东、西北、朝鲜或西南军务的本年调兵与战报记录跨区域动员，补给与兵额需要独立核验。（材料核查重点：诏令或奏议的形成与发受链）”发生于明。本条把背景说明为前期边防、地方武装或跨区域资源调动的压力推动了该行动。，并以它改变了后续调兵、补给或地方秩序的条件；战果和伤亡须分开核对。概括可见影响。要理解它的社会后果，不能止于宫廷或战场：土地、赋役、司法、道路、性别化家庭劳动、城市市场、书院寺观和边海网络都会影响地方如何接收或改写制度。

本条依据《神宗实录》万历四年条；《明史·兵志》及相关列传；边镇、地方志或对方材料。这些来源能够较稳妥地支持所记年序、诏令、任命或特定地点的行动，但无法自动证明全国的财产、人口、识字、信仰或动机。账本保留的证据说明是“桥接条目只标示可回查的年度问题与材料链；实录有中央选择性，崇祯期尤其需要奏疏、地方、朝鲜及满文材料交叉。；本条特别核对诏令或奏议的形成与发受链。”。所以，官府标签、奏疏数字与后来的道德评语必须放回其文书目的，不能被写成不经分区核验的社会全貌。

这条事件更适合作为一组关系变化的锚点：中央方案需由地方官、书吏、士绅、宗族、商人与劳动者共同转译；不同家庭面对的税役、婚姻、迁徙和安全风险并不相等。材料只让我们看见其中有限的记录者。承认可见与不可见的界线，才能使制度史、文化史和区域史互相校正。

在账本条目\texttt{undefined}中，1576（万历四年）的“考成、人事、立储和留中等本年材料标记皇帝、内阁与言官的协调关系，不可只用党争标签解释。”发生于明。本条把背景说明为继承、官僚协调或皇权运作的压力使问题进入朝廷程序。，并以它影响后续人事与政策议程；动机和派系标签不能由单一叙事裁定。概括可见影响。要理解它的社会后果，不能止于宫廷或战场：土地、赋役、司法、道路、性别化家庭劳动、城市市场、书院寺观和边海网络都会影响地方如何接收或改写制度。

本条依据《神宗实录》万历四年条；《明史·本纪》及相关列传；诏令、奏疏或会典版本。这些来源能够较稳妥地支持所记年序、诏令、任命或特定地点的行动，但无法自动证明全国的财产、人口、识字、信仰或动机。账本保留的证据说明是“桥接条目只标示可回查的年度问题与材料链；实录有中央选择性，崇祯期尤其需要奏疏、地方、朝鲜及满文材料交叉。”。所以，官府标签、奏疏数字与后来的道德评语必须放回其文书目的，不能被写成不经分区核验的社会全貌。

这条事件更适合作为一组关系变化的锚点：中央方案需由地方官、书吏、士绅、宗族、商人与劳动者共同转译；不同家庭面对的税役、婚姻、迁徙和安全风险并不相等。材料只让我们看见其中有限的记录者。承认可见与不可见的界线，才能使制度史、文化史和区域史互相校正。

在账本条目\texttt{undefined}中，1576（万历四年）的“辽东、西北、朝鲜或西南军务的本年调兵与战报记录跨区域动员，补给与兵额需要独立核验。（材料核查重点：报称信息的来源、抄传与行政分类）”发生于明。本条把背景说明为前期边防、地方武装或跨区域资源调动的压力推动了该行动。，并以它改变了后续调兵、补给或地方秩序的条件；战果和伤亡须分开核对。概括可见影响。要理解它的社会后果，不能止于宫廷或战场：土地、赋役、司法、道路、性别化家庭劳动、城市市场、书院寺观和边海网络都会影响地方如何接收或改写制度。

本条依据《神宗实录》万历四年条；《明史·兵志》及相关列传；边镇、地方志或对方材料。这些来源能够较稳妥地支持所记年序、诏令、任命或特定地点的行动，但无法自动证明全国的财产、人口、识字、信仰或动机。账本保留的证据说明是“桥接条目只标示可回查的年度问题与材料链；实录有中央选择性，崇祯期尤其需要奏疏、地方、朝鲜及满文材料交叉。；本条特别核对报称信息的来源、抄传与行政分类。”。所以，官府标签、奏疏数字与后来的道德评语必须放回其文书目的，不能被写成不经分区核验的社会全貌。

这条事件更适合作为一组关系变化的锚点：中央方案需由地方官、书吏、士绅、宗族、商人与劳动者共同转译；不同家庭面对的税役、婚姻、迁徙和安全风险并不相等。材料只让我们看见其中有限的记录者。承认可见与不可见的界线，才能使制度史、文化史和区域史互相校正。

\subsection{城市、文化与知识流通}

在账本条目\texttt{undefined}中，1576（万历四年）的“本年灾荒、河工和地方赈济材料可用于研究风险分配，灾害叙事和统计数字应区分。”发生于明。本条把背景说明为自然风险与地方报告促成朝廷或地方的应对记录。，并以赈济、迁徙和损失的实际范围仍须以受灾地区材料分别核验。概括可见影响。要理解它的社会后果，不能止于宫廷或战场：土地、赋役、司法、道路、性别化家庭劳动、城市市场、书院寺观和边海网络都会影响地方如何接收或改写制度。

本条依据《神宗实录》万历四年条；《明史·五行志》；受灾府县地方志与奏疏。这些来源能够较稳妥地支持所记年序、诏令、任命或特定地点的行动，但无法自动证明全国的财产、人口、识字、信仰或动机。账本保留的证据说明是“桥接条目只标示可回查的年度问题与材料链；实录有中央选择性，崇祯期尤其需要奏疏、地方、朝鲜及满文材料交叉。”。所以，官府标签、奏疏数字与后来的道德评语必须放回其文书目的，不能被写成不经分区核验的社会全貌。

这条事件更适合作为一组关系变化的锚点：中央方案需由地方官、书吏、士绅、宗族、商人与劳动者共同转译；不同家庭面对的税役、婚姻、迁徙和安全风险并不相等。材料只让我们看见其中有限的记录者。承认可见与不可见的界线，才能使制度史、文化史和区域史互相校正。

在账本条目\texttt{undefined}中，1576（万历四年）的“朝鲜战争、海上贸易和朝贡使节的本年多方记录提供跨政权比较，外交语言不能替代地方经验。”发生于明。本条把背景说明为朝贡名义、边境安全与港口贸易的利益使交涉持续发生。，并以它为后续册封、互市或海防安排留下文书与跨政权比较线索。概括可见影响。要理解它的社会后果，不能止于宫廷或战场：土地、赋役、司法、道路、性别化家庭劳动、城市市场、书院寺观和边海网络都会影响地方如何接收或改写制度。

本条依据《神宗实录》万历四年条；《明史·外国传》；《朝鲜王朝实录》、琉球《历代宝案》或海外档案。这些来源能够较稳妥地支持所记年序、诏令、任命或特定地点的行动，但无法自动证明全国的财产、人口、识字、信仰或动机。账本保留的证据说明是“桥接条目只标示可回查的年度问题与材料链；实录有中央选择性，崇祯期尤其需要奏疏、地方、朝鲜及满文材料交叉。”。所以，官府标签、奏疏数字与后来的道德评语必须放回其文书目的，不能被写成不经分区核验的社会全貌。

这条事件更适合作为一组关系变化的锚点：中央方案需由地方官、书吏、士绅、宗族、商人与劳动者共同转译；不同家庭面对的税役、婚姻、迁徙和安全风险并不相等。材料只让我们看见其中有限的记录者。承认可见与不可见的界线，才能使制度史、文化史和区域史互相校正。

在账本条目\texttt{undefined}中，1577（万历五年）的“辽东、西北、朝鲜或西南军务的本年调兵与战报记录跨区域动员，补给与兵额需要独立核验。”发生于明。本条把背景说明为前期边防、地方武装或跨区域资源调动的压力推动了该行动。，并以它改变了后续调兵、补给或地方秩序的条件；战果和伤亡须分开核对。概括可见影响。要理解它的社会后果，不能止于宫廷或战场：土地、赋役、司法、道路、性别化家庭劳动、城市市场、书院寺观和边海网络都会影响地方如何接收或改写制度。

本条依据《神宗实录》万历五年条；《明史·兵志》及相关列传；边镇、地方志或对方材料。这些来源能够较稳妥地支持所记年序、诏令、任命或特定地点的行动，但无法自动证明全国的财产、人口、识字、信仰或动机。账本保留的证据说明是“桥接条目只标示可回查的年度问题与材料链；实录有中央选择性，崇祯期尤其需要奏疏、地方、朝鲜及满文材料交叉。”。所以，官府标签、奏疏数字与后来的道德评语必须放回其文书目的，不能被写成不经分区核验的社会全貌。

这条事件更适合作为一组关系变化的锚点：中央方案需由地方官、书吏、士绅、宗族、商人与劳动者共同转译；不同家庭面对的税役、婚姻、迁徙和安全风险并不相等。材料只让我们看见其中有限的记录者。承认可见与不可见的界线，才能使制度史、文化史和区域史互相校正。

在账本条目\texttt{undefined}中，1577（万历五年）的“朝鲜战争、海上贸易和朝贡使节的本年多方记录提供跨政权比较，外交语言不能替代地方经验。”发生于明。本条把背景说明为朝贡名义、边境安全与港口贸易的利益使交涉持续发生。，并以它为后续册封、互市或海防安排留下文书与跨政权比较线索。概括可见影响。要理解它的社会后果，不能止于宫廷或战场：土地、赋役、司法、道路、性别化家庭劳动、城市市场、书院寺观和边海网络都会影响地方如何接收或改写制度。

本条依据《神宗实录》万历五年条；《明史·外国传》；《朝鲜王朝实录》、琉球《历代宝案》或海外档案。这些来源能够较稳妥地支持所记年序、诏令、任命或特定地点的行动，但无法自动证明全国的财产、人口、识字、信仰或动机。账本保留的证据说明是“桥接条目只标示可回查的年度问题与材料链；实录有中央选择性，崇祯期尤其需要奏疏、地方、朝鲜及满文材料交叉。”。所以，官府标签、奏疏数字与后来的道德评语必须放回其文书目的，不能被写成不经分区核验的社会全貌。

这条事件更适合作为一组关系变化的锚点：中央方案需由地方官、书吏、士绅、宗族、商人与劳动者共同转译；不同家庭面对的税役、婚姻、迁徙和安全风险并不相等。材料只让我们看见其中有限的记录者。承认可见与不可见的界线，才能使制度史、文化史和区域史互相校正。

在账本条目\texttt{undefined}中，1577（万历五年）的“清丈、一条鞭、漕运或军饷的本年文书构成万历财政的可核查节点，定额、报数和实收必须分列。”发生于明。本条把背景说明为财政征解、交通工程或货币支付的压力使相关措施进入政务。，并以其法定安排不等于地方执行，后续须以地域材料追踪负担与成效。概括可见影响。要理解它的社会后果，不能止于宫廷或战场：土地、赋役、司法、道路、性别化家庭劳动、城市市场、书院寺观和边海网络都会影响地方如何接收或改写制度。

本条依据《神宗实录》万历五年条；《明会典》相关条目；地方志、黄册/契约/河工材料。这些来源能够较稳妥地支持所记年序、诏令、任命或特定地点的行动，但无法自动证明全国的财产、人口、识字、信仰或动机。账本保留的证据说明是“桥接条目只标示可回查的年度问题与材料链；实录有中央选择性，崇祯期尤其需要奏疏、地方、朝鲜及满文材料交叉。”。所以，官府标签、奏疏数字与后来的道德评语必须放回其文书目的，不能被写成不经分区核验的社会全貌。

这条事件更适合作为一组关系变化的锚点：中央方案需由地方官、书吏、士绅、宗族、商人与劳动者共同转译；不同家庭面对的税役、婚姻、迁徙和安全风险并不相等。材料只让我们看见其中有限的记录者。承认可见与不可见的界线，才能使制度史、文化史和区域史互相校正。

在账本条目\texttt{undefined}中，1577（万历五年）的“书院、科举、出版与宗教活动的本年线索可连接知识网络，存世文本有阶层与地域偏差。”发生于明。本条把背景说明为制度、教育或知识传播的需要推动了相关行动或文本形成。，并以它提供制度和传播史的锚点，不能据此推断社会整体接受。概括可见影响。要理解它的社会后果，不能止于宫廷或战场：土地、赋役、司法、道路、性别化家庭劳动、城市市场、书院寺观和边海网络都会影响地方如何接收或改写制度。

本条依据《神宗实录》万历五年条；《明史·选举志》或《艺文志》；文集、碑刻或书目材料。这些来源能够较稳妥地支持所记年序、诏令、任命或特定地点的行动，但无法自动证明全国的财产、人口、识字、信仰或动机。账本保留的证据说明是“桥接条目只标示可回查的年度问题与材料链；实录有中央选择性，崇祯期尤其需要奏疏、地方、朝鲜及满文材料交叉。”。所以，官府标签、奏疏数字与后来的道德评语必须放回其文书目的，不能被写成不经分区核验的社会全貌。

这条事件更适合作为一组关系变化的锚点：中央方案需由地方官、书吏、士绅、宗族、商人与劳动者共同转译；不同家庭面对的税役、婚姻、迁徙和安全风险并不相等。材料只让我们看见其中有限的记录者。承认可见与不可见的界线，才能使制度史、文化史和区域史互相校正。

在账本条目\texttt{undefined}中，1577（万历五年）的“潘季驯受命主持河防，提出束水攻沙并组织黄河堤防与运河保运工程”发生于明与黄河、运河沿线。本条把背景说明为黄河决溢威胁运河漕运、沿河农田和漕粮北运，中央需在河工经费、役夫征发与水势变化之间选择方案。，并以堤防与导流工程加强了对若干河段的控制并保障部分漕运，亦持续要求地方承担经费和劳役；河患不因一次治理而消除。概括可见影响。要理解它的社会后果，不能止于宫廷或战场：土地、赋役、司法、道路、性别化家庭劳动、城市市场、书院寺观和边海网络都会影响地方如何接收或改写制度。

本条依据《神宗实录》万历五年河道条；潘季驯《河防一览》；《明史·河渠志》。这些来源能够较稳妥地支持所记年序、诏令、任命或特定地点的行动，但无法自动证明全国的财产、人口、识字、信仰或动机。账本保留的证据说明是“任命、奏议和工程命令可靠；治河成效随河段、水情和维护周期而异，后出技术赞誉不能证明所有堤段安全或漕运畅通。”。所以，官府标签、奏疏数字与后来的道德评语必须放回其文书目的，不能被写成不经分区核验的社会全貌。

这条事件更适合作为一组关系变化的锚点：中央方案需由地方官、书吏、士绅、宗族、商人与劳动者共同转译；不同家庭面对的税役、婚姻、迁徙和安全风险并不相等。材料只让我们看见其中有限的记录者。承认可见与不可见的界线，才能使制度史、文化史和区域史互相校正。

在账本条目\texttt{undefined}中，1577（万历五年）的“张居正父丧期间奉旨夺情留任，引发朝臣围绕礼制与政务的争论”发生于明。本条把背景说明为张居正主持的改革需要连续推动，太后与皇帝担忧其离职；传统士大夫伦理则强调居丧服制。，并以张居正得以留任，但政治反对扩大；礼制争论成为评价其权力边界和万历初政的重要线索。概括可见影响。要理解它的社会后果，不能止于宫廷或战场：土地、赋役、司法、道路、性别化家庭劳动、城市市场、书院寺观和边海网络都会影响地方如何接收或改写制度。

本条依据《神宗实录》万历五年条；《明史·张居正传》；张居正、吴中行等奏疏。这些来源能够较稳妥地支持所记年序、诏令、任命或特定地点的行动，但无法自动证明全国的财产、人口、识字、信仰或动机。账本保留的证据说明是“夺情诏令和争论可靠；反对者并非只出于礼制或派系，具体立场应回到各人奏疏。”。所以，官府标签、奏疏数字与后来的道德评语必须放回其文书目的，不能被写成不经分区核验的社会全貌。

这条事件更适合作为一组关系变化的锚点：中央方案需由地方官、书吏、士绅、宗族、商人与劳动者共同转译；不同家庭面对的税役、婚姻、迁徙和安全风险并不相等。材料只让我们看见其中有限的记录者。承认可见与不可见的界线，才能使制度史、文化史和区域史互相校正。

在账本条目\texttt{undefined}中，1578（万历六年）的“葡萄牙人获准前往广州”发生于明与海上、朝贡相关政权。本条把背景说明为朝贡名义、边境安全与港口贸易的利益使交涉持续发生。，并以它为后续册封、互市或海防安排留下文书与跨政权比较线索。概括可见影响。要理解它的社会后果，不能止于宫廷或战场：土地、赋役、司法、道路、性别化家庭劳动、城市市场、书院寺观和边海网络都会影响地方如何接收或改写制度。

本条依据《神宗实录》万历六年条；《明史·外国传》；《朝鲜王朝实录》、琉球《历代宝案》或海外档案。这些来源能够较稳妥地支持所记年序、诏令、任命或特定地点的行动，但无法自动证明全国的财产、人口、识字、信仰或动机。账本保留的证据说明是“译写候选以编年材料定位；实录记载反映朝廷选择，崇祯期须另以奏疏、地方及敌对政权材料交叉。”。所以，官府标签、奏疏数字与后来的道德评语必须放回其文书目的，不能被写成不经分区核验的社会全貌。

这条事件更适合作为一组关系变化的锚点：中央方案需由地方官、书吏、士绅、宗族、商人与劳动者共同转译；不同家庭面对的税役、婚姻、迁徙和安全风险并不相等。材料只让我们看见其中有限的记录者。承认可见与不可见的界线，才能使制度史、文化史和区域史互相校正。

在账本条目\texttt{undefined}中，1578（万历六年）的“书院、科举、出版与宗教活动的本年线索可连接知识网络，存世文本有阶层与地域偏差。”发生于明。本条把背景说明为制度、教育或知识传播的需要推动了相关行动或文本形成。，并以它提供制度和传播史的锚点，不能据此推断社会整体接受。概括可见影响。要理解它的社会后果，不能止于宫廷或战场：土地、赋役、司法、道路、性别化家庭劳动、城市市场、书院寺观和边海网络都会影响地方如何接收或改写制度。

本条依据《神宗实录》万历六年条；《明史·选举志》或《艺文志》；文集、碑刻或书目材料。这些来源能够较稳妥地支持所记年序、诏令、任命或特定地点的行动，但无法自动证明全国的财产、人口、识字、信仰或动机。账本保留的证据说明是“桥接条目只标示可回查的年度问题与材料链；实录有中央选择性，崇祯期尤其需要奏疏、地方、朝鲜及满文材料交叉。”。所以，官府标签、奏疏数字与后来的道德评语必须放回其文书目的，不能被写成不经分区核验的社会全貌。

这条事件更适合作为一组关系变化的锚点：中央方案需由地方官、书吏、士绅、宗族、商人与劳动者共同转译；不同家庭面对的税役、婚姻、迁徙和安全风险并不相等。材料只让我们看见其中有限的记录者。承认可见与不可见的界线，才能使制度史、文化史和区域史互相校正。

在账本条目\texttt{undefined}中，1578（万历六年）的“考成、人事、立储和留中等本年材料标记皇帝、内阁与言官的协调关系，不可只用党争标签解释。”发生于明。本条把背景说明为继承、官僚协调或皇权运作的压力使问题进入朝廷程序。，并以它影响后续人事与政策议程；动机和派系标签不能由单一叙事裁定。概括可见影响。要理解它的社会后果，不能止于宫廷或战场：土地、赋役、司法、道路、性别化家庭劳动、城市市场、书院寺观和边海网络都会影响地方如何接收或改写制度。

本条依据《神宗实录》万历六年条；《明史·本纪》及相关列传；诏令、奏疏或会典版本。这些来源能够较稳妥地支持所记年序、诏令、任命或特定地点的行动，但无法自动证明全国的财产、人口、识字、信仰或动机。账本保留的证据说明是“桥接条目只标示可回查的年度问题与材料链；实录有中央选择性，崇祯期尤其需要奏疏、地方、朝鲜及满文材料交叉。”。所以，官府标签、奏疏数字与后来的道德评语必须放回其文书目的，不能被写成不经分区核验的社会全貌。

这条事件更适合作为一组关系变化的锚点：中央方案需由地方官、书吏、士绅、宗族、商人与劳动者共同转译；不同家庭面对的税役、婚姻、迁徙和安全风险并不相等。材料只让我们看见其中有限的记录者。承认可见与不可见的界线，才能使制度史、文化史和区域史互相校正。

\subsection{环境、边防与海外连接}

在账本条目\texttt{undefined}中，1578（万历六年）的“本年灾荒、河工和地方赈济材料可用于研究风险分配，灾害叙事和统计数字应区分。”发生于明。本条把背景说明为自然风险与地方报告促成朝廷或地方的应对记录。，并以赈济、迁徙和损失的实际范围仍须以受灾地区材料分别核验。概括可见影响。要理解它的社会后果，不能止于宫廷或战场：土地、赋役、司法、道路、性别化家庭劳动、城市市场、书院寺观和边海网络都会影响地方如何接收或改写制度。

本条依据《神宗实录》万历六年条；《明史·五行志》；受灾府县地方志与奏疏。这些来源能够较稳妥地支持所记年序、诏令、任命或特定地点的行动，但无法自动证明全国的财产、人口、识字、信仰或动机。账本保留的证据说明是“桥接条目只标示可回查的年度问题与材料链；实录有中央选择性，崇祯期尤其需要奏疏、地方、朝鲜及满文材料交叉。”。所以，官府标签、奏疏数字与后来的道德评语必须放回其文书目的，不能被写成不经分区核验的社会全貌。

这条事件更适合作为一组关系变化的锚点：中央方案需由地方官、书吏、士绅、宗族、商人与劳动者共同转译；不同家庭面对的税役、婚姻、迁徙和安全风险并不相等。材料只让我们看见其中有限的记录者。承认可见与不可见的界线，才能使制度史、文化史和区域史互相校正。

在账本条目\texttt{undefined}中，1578（万历六年）六月的“万历帝册立王氏为皇后，完成幼帝成年后的皇后册立仪式”发生于明。本条把背景说明为万历帝已亲政，皇室需要依礼确定皇后以组织后宫名分、祭祀和未来继嗣的程序。，并以王氏取得皇后名号，后宫礼制与外戚关系有了正式支点；继承问题仍取决于子嗣、册立和宫廷政治。概括可见影响。要理解它的社会后果，不能止于宫廷或战场：土地、赋役、司法、道路、性别化家庭劳动、城市市场、书院寺观和边海网络都会影响地方如何接收或改写制度。

本条依据《神宗实录》万历六年六月乙卯条；《明史·孝端显皇后传》；明代后妃册立仪文。这些来源能够较稳妥地支持所记年序、诏令、任命或特定地点的行动，但无法自动证明全国的财产、人口、识字、信仰或动机。账本保留的证据说明是“册立日期、礼仪和诏令可靠；礼制记录突出朝廷秩序，不能据此推知帝后关系、后宫实际权力或后继生育结果。”。所以，官府标签、奏疏数字与后来的道德评语必须放回其文书目的，不能被写成不经分区核验的社会全貌。

这条事件更适合作为一组关系变化的锚点：中央方案需由地方官、书吏、士绅、宗族、商人与劳动者共同转译；不同家庭面对的税役、婚姻、迁徙和安全风险并不相等。材料只让我们看见其中有限的记录者。承认可见与不可见的界线，才能使制度史、文化史和区域史互相校正。

在账本条目\texttt{undefined}中，1578（万历六年）的“朝廷令各省清丈田亩、核查隐漏，为赋役重整收集册籍”发生于明与各省田土社会。本条把背景说明为旧有登记与实际占有、赋额脱节，中央财政和改革政策需要更可比的田亩、户役信息。，并以各地重编册籍并暴露隐漏、产权和役额争议；清丈为一条鞭的区域推行提供条件，却不决定其统一形式。概括可见影响。要理解它的社会后果，不能止于宫廷或战场：土地、赋役、司法、道路、性别化家庭劳动、城市市场、书院寺观和边海网络都会影响地方如何接收或改写制度。

本条依据《神宗实录》万历六年户部条；《明史·食货志》；徽州、江南等地清丈册与契约。这些来源能够较稳妥地支持所记年序、诏令、任命或特定地点的行动，但无法自动证明全国的财产、人口、识字、信仰或动机。账本保留的证据说明是“清丈命令可靠；各省启动、完结与报亩年份不同，中央汇总数不能当作同一年度的全国实测。”。所以，官府标签、奏疏数字与后来的道德评语必须放回其文书目的，不能被写成不经分区核验的社会全貌。

这条事件更适合作为一组关系变化的锚点：中央方案需由地方官、书吏、士绅、宗族、商人与劳动者共同转译；不同家庭面对的税役、婚姻、迁徙和安全风险并不相等。材料只让我们看见其中有限的记录者。承认可见与不可见的界线，才能使制度史、文化史和区域史互相校正。

在账本条目\texttt{undefined}中，1579（万历七年）的“东林运动：私立东林书院全部关闭”发生于明。本条把背景说明为继承、官僚协调或皇权运作的压力使问题进入朝廷程序。，并以它影响后续人事与政策议程；动机和派系标签不能由单一叙事裁定。概括可见影响。要理解它的社会后果，不能止于宫廷或战场：土地、赋役、司法、道路、性别化家庭劳动、城市市场、书院寺观和边海网络都会影响地方如何接收或改写制度。

本条依据《神宗实录》万历七年条；《明史·本纪》及相关列传；诏令、奏疏或会典版本。这些来源能够较稳妥地支持所记年序、诏令、任命或特定地点的行动，但无法自动证明全国的财产、人口、识字、信仰或动机。账本保留的证据说明是“译写候选以编年材料定位；实录记载反映朝廷选择，崇祯期须另以奏疏、地方及敌对政权材料交叉。”。所以，官府标签、奏疏数字与后来的道德评语必须放回其文书目的，不能被写成不经分区核验的社会全貌。

这条事件更适合作为一组关系变化的锚点：中央方案需由地方官、书吏、士绅、宗族、商人与劳动者共同转译；不同家庭面对的税役、婚姻、迁徙和安全风险并不相等。材料只让我们看见其中有限的记录者。承认可见与不可见的界线，才能使制度史、文化史和区域史互相校正。

在账本条目\texttt{undefined}中，1579（万历七年）的“朝鲜战争、海上贸易和朝贡使节的本年多方记录提供跨政权比较，外交语言不能替代地方经验。”发生于明。本条把背景说明为朝贡名义、边境安全与港口贸易的利益使交涉持续发生。，并以它为后续册封、互市或海防安排留下文书与跨政权比较线索。概括可见影响。要理解它的社会后果，不能止于宫廷或战场：土地、赋役、司法、道路、性别化家庭劳动、城市市场、书院寺观和边海网络都会影响地方如何接收或改写制度。

本条依据《神宗实录》万历七年条；《明史·外国传》；《朝鲜王朝实录》、琉球《历代宝案》或海外档案。这些来源能够较稳妥地支持所记年序、诏令、任命或特定地点的行动，但无法自动证明全国的财产、人口、识字、信仰或动机。账本保留的证据说明是“桥接条目只标示可回查的年度问题与材料链；实录有中央选择性，崇祯期尤其需要奏疏、地方、朝鲜及满文材料交叉。”。所以，官府标签、奏疏数字与后来的道德评语必须放回其文书目的，不能被写成不经分区核验的社会全貌。

这条事件更适合作为一组关系变化的锚点：中央方案需由地方官、书吏、士绅、宗族、商人与劳动者共同转译；不同家庭面对的税役、婚姻、迁徙和安全风险并不相等。材料只让我们看见其中有限的记录者。承认可见与不可见的界线，才能使制度史、文化史和区域史互相校正。

在账本条目\texttt{undefined}中，1579（万历七年）的“本年灾荒、河工和地方赈济材料可用于研究风险分配，灾害叙事和统计数字应区分。”发生于明。本条把背景说明为自然风险与地方报告促成朝廷或地方的应对记录。，并以赈济、迁徙和损失的实际范围仍须以受灾地区材料分别核验。概括可见影响。要理解它的社会后果，不能止于宫廷或战场：土地、赋役、司法、道路、性别化家庭劳动、城市市场、书院寺观和边海网络都会影响地方如何接收或改写制度。

本条依据《神宗实录》万历七年条；《明史·五行志》；受灾府县地方志与奏疏。这些来源能够较稳妥地支持所记年序、诏令、任命或特定地点的行动，但无法自动证明全国的财产、人口、识字、信仰或动机。账本保留的证据说明是“桥接条目只标示可回查的年度问题与材料链；实录有中央选择性，崇祯期尤其需要奏疏、地方、朝鲜及满文材料交叉。”。所以，官府标签、奏疏数字与后来的道德评语必须放回其文书目的，不能被写成不经分区核验的社会全貌。

这条事件更适合作为一组关系变化的锚点：中央方案需由地方官、书吏、士绅、宗族、商人与劳动者共同转译；不同家庭面对的税役、婚姻、迁徙和安全风险并不相等。材料只让我们看见其中有限的记录者。承认可见与不可见的界线，才能使制度史、文化史和区域史互相校正。

在账本条目\texttt{undefined}中，1579（万历七年）的“辽东、西北、朝鲜或西南军务的本年调兵与战报记录跨区域动员，补给与兵额需要独立核验。”发生于明。本条把背景说明为前期边防、地方武装或跨区域资源调动的压力推动了该行动。，并以它改变了后续调兵、补给或地方秩序的条件；战果和伤亡须分开核对。概括可见影响。要理解它的社会后果，不能止于宫廷或战场：土地、赋役、司法、道路、性别化家庭劳动、城市市场、书院寺观和边海网络都会影响地方如何接收或改写制度。

本条依据《神宗实录》万历七年条；《明史·兵志》及相关列传；边镇、地方志或对方材料。这些来源能够较稳妥地支持所记年序、诏令、任命或特定地点的行动，但无法自动证明全国的财产、人口、识字、信仰或动机。账本保留的证据说明是“桥接条目只标示可回查的年度问题与材料链；实录有中央选择性，崇祯期尤其需要奏疏、地方、朝鲜及满文材料交叉。”。所以，官府标签、奏疏数字与后来的道德评语必须放回其文书目的，不能被写成不经分区核验的社会全貌。

这条事件更适合作为一组关系变化的锚点：中央方案需由地方官、书吏、士绅、宗族、商人与劳动者共同转译；不同家庭面对的税役、婚姻、迁徙和安全风险并不相等。材料只让我们看见其中有限的记录者。承认可见与不可见的界线，才能使制度史、文化史和区域史互相校正。

在账本条目\texttt{undefined}中，1579（万历七年）的“清丈、一条鞭、漕运或军饷的本年文书构成万历财政的可核查节点，定额、报数和实收必须分列。”发生于明。本条把背景说明为财政征解、交通工程或货币支付的压力使相关措施进入政务。，并以其法定安排不等于地方执行，后续须以地域材料追踪负担与成效。概括可见影响。要理解它的社会后果，不能止于宫廷或战场：土地、赋役、司法、道路、性别化家庭劳动、城市市场、书院寺观和边海网络都会影响地方如何接收或改写制度。

本条依据《神宗实录》万历七年条；《明会典》相关条目；地方志、黄册/契约/河工材料。这些来源能够较稳妥地支持所记年序、诏令、任命或特定地点的行动，但无法自动证明全国的财产、人口、识字、信仰或动机。账本保留的证据说明是“桥接条目只标示可回查的年度问题与材料链；实录有中央选择性，崇祯期尤其需要奏疏、地方、朝鲜及满文材料交叉。”。所以，官府标签、奏疏数字与后来的道德评语必须放回其文书目的，不能被写成不经分区核验的社会全貌。

这条事件更适合作为一组关系变化的锚点：中央方案需由地方官、书吏、士绅、宗族、商人与劳动者共同转译；不同家庭面对的税役、婚姻、迁徙和安全风险并不相等。材料只让我们看见其中有限的记录者。承认可见与不可见的界线，才能使制度史、文化史和区域史互相校正。

在账本条目\texttt{undefined}中，1579（万历七年）的“张居正政府以禁讲学、限制书院活动的诏令约束地方士人结社”发生于明与各地书院、士人。本条把背景说明为张居正担忧讲学社群与地方结交影响官员考成和政治批评，中央希望将士人活动重新纳入官学与官僚秩序。，并以部分书院公开活动受限，士人转向家塾、文会或其他交往形式；禁令也为后来的学术自治记忆提供素材。概括可见影响。要理解它的社会后果，不能止于宫廷或战场：土地、赋役、司法、道路、性别化家庭劳动、城市市场、书院寺观和边海网络都会影响地方如何接收或改写制度。

本条依据《神宗实录》万历七年书院条；张居正公牍与奏疏；书院碑记及地方志。这些来源能够较稳妥地支持所记年序、诏令、任命或特定地点的行动，但无法自动证明全国的财产、人口、识字、信仰或动机。账本保留的证据说明是“禁令和中央政策方向可靠；各地书院是否停讲、改名或私下延续须以地方志、碑刻和文集逐处核验，不能以诏令写成思想生活完全中断。”。所以，官府标签、奏疏数字与后来的道德评语必须放回其文书目的，不能被写成不经分区核验的社会全貌。

这条事件更适合作为一组关系变化的锚点：中央方案需由地方官、书吏、士绅、宗族、商人与劳动者共同转译；不同家庭面对的税役、婚姻、迁徙和安全风险并不相等。材料只让我们看见其中有限的记录者。承认可见与不可见的界线，才能使制度史、文化史和区域史互相校正。
